% Section body only — the builder emits \section{...}\label{sec:framework}.
% Primary sources: docs/derivations/G2c_gray_a9_a10_mapping.md,
%                  docs/derivations/G2a_completion_sky_marginal_4pi.md,
%                  docs/derivations/G2b_host_z_volume_prior.md,
%                  derivations/dark_siren_likelihood.md
The estimator analysed in this paper is the galaxy-catalogue
(``dark-siren'') Hubble-constant likelihood of \citet{Gray:2019ksv},
as implemented for a simulated LISA extreme-mass-ratio-inspiral (EMRI)
campaign \citep{Laghi:2021pqk} evaluated against the GLADE+ catalogue
\citep{Dalya:2021ewn}.  This section assembles the
complete single-event likelihood and fixes the notation used
throughout the paper; Section~\ref{sec:pitfall} then isolates the
normalization choices hidden inside it, and
Appendix~\ref{sec:app:gray} gives the term-by-term mapping onto the
published equations of \citet{Gray:2019ksv}.

\subsection{Notation, cosmology, and single-event data}
\label{sec:framework:notation}

We write the Hubble constant through the dimensionless parameter
$h \equiv H_0 / (100~\mathrm{km\,s^{-1}\,Mpc^{-1}})$.  All redshifts
$z$ are CMB-frame redshifts, and $\sigma_z$ denotes the (per-galaxy)
photometric redshift uncertainty.  We assume a flat $\Lambda$CDM
background, in which the luminosity distance is
\begin{equation}
  d_L(z, h) \;=\; \frac{c\,(1+z)}{H_0} \int_0^z
  \frac{\mathrm{d}z'}{E(z')},
  \qquad
  E(z) = \sqrt{\Omega_\mathrm{m}(1+z)^3 + \Omega_\Lambda},
  \label{eq:dl}
\end{equation}
with $\Omega_\Lambda = 1 - \Omega_\mathrm{m}$ \citep{Hogg:1999ad}.
Because the data are drawn from a simulated EMRI campaign (on the real
GLADE+ galaxy field), all statements about recovering $h$ refer to
estimator calibration against the injected simulation truth,
$h = 0.73$,
% src: master_thesis_code/constants.py (H = 0.73, fiducial simulation value)
never to a measurement of the actual expansion rate.

Each detected event $i$ contributes data $x_i$ and a detection
indicator $D_i$ (the event's signal-to-noise ratio exceeds the
detection threshold $\rho_\mathrm{thr} = 20$).
% src: master_thesis_code/constants.py (SNR_THRESHOLD = 20)
With a prior $\pi(h)$ (uniform here), the posterior over $h$ from
$N_\mathrm{det}$ events is
\begin{equation}
  p\big(h \mid \{x_i\}\big) \;\propto\; \pi(h)
  \prod_{i=1}^{N_\mathrm{det}} p_i(h),
  \qquad
  p_i(h) \equiv p(x_i \mid D_i, h).
  \label{eq:posterior}
\end{equation}

The per-event measurement is summarized by a Gaussian likelihood
obtained from the Cram\'er--Rao bound of a 14-parameter EMRI Fisher
analysis
% src: master_thesis_code/datamodels/parameter_space.py (14 Parameter fields)
\citep{Vallisneri:2007ev}, projected onto the three parameters that
carry cosmological information: the ecliptic sky coordinates
$(\phi, \theta)$ and the luminosity-distance fraction
$u \equiv d_L / \hat d_L$, where $\hat d_L$ is the event's
maximum-likelihood distance.  Writing $\hat\mu_i = (\hat\phi_i,
\hat\theta_i, 1)$ for the maximum-likelihood point and $\Sigma_i$ for
the projected covariance,
\begin{equation}
  p\big(x_i \mid z, \Omega, h\big) \;\simeq\;
  \mathcal{N}_3\!\big( (\phi, \theta, u(z,h));\; \hat\mu_i,\, \Sigma_i \big),
  \qquad
  u(z,h) = \frac{d_L(z,h)}{\hat d_L},
  \label{eq:gwlike}
\end{equation}
% src: docs/derivations/G2a_completion_sky_marginal_4pi.md §1
where $\Omega = (\phi,\theta)$ denotes sky position.  Equation
\eqref{eq:gwlike} is how $h$ enters the inference: for each trial $h$,
the map \eqref{eq:dl} converts a candidate redshift into a distance
fraction at which the Gaussian is evaluated.  One convention matters
for everything that follows: the Fisher matrix is computed with
respect to the bare coordinates $(\phi,\theta)$, so
\eqref{eq:gwlike} is a probability density with respect to
$\mathrm{d}\phi\,\mathrm{d}\theta\,\mathrm{d}u$, \emph{not} with
respect to the solid-angle measure
$\mathrm{d}\Omega = \sin\theta\,\mathrm{d}\theta\,\mathrm{d}\phi$;
sky integrals against isotropic priors must therefore carry the
$\sin\theta$ Jacobian explicitly (Appendix~\ref{sec:app:sky}).

\subsection{Hypothesis decomposition}
\label{sec:framework:decomposition}

Following \citet{Gray:2019ksv}, the likelihood marginalizes over the
two exhaustive hypotheses $G$ (the true host galaxy is contained in
the flux-limited catalogue) and $\bar G$ (it is not):
\begin{equation}
  p_i(h) \;=\;
  p(x_i \mid G, D_i, h)\, p(G \mid D_i, h)
  \;+\;
  p(x_i \mid \bar G, D_i, h)\, p(\bar G \mid D_i, h).
  \label{eq:mixture}
\end{equation}
% src: docs/derivations/G2c_gray_a9_a10_mapping.md §0 (Gray main-text Eq. 9)
We refer to the second term as the \emph{completion term}: it restores
the contribution of hosts the catalogue is missing.  The four
ingredients of \eqref{eq:mixture} --- the in-catalogue likelihood, the
completion term, and the two hypothesis weights --- are constructed in
the next three subsections and reassembled in
equation~\eqref{eq:assembled}.

\subsection{In-catalogue term}
\label{sec:framework:incat}

The in-catalogue likelihood is a rate-weighted ratio of sums over
candidate host galaxies $g$,
\begin{equation}
  L_\mathrm{cat}(h) \;\equiv\; p(x_i \mid G, D_i, h)
  \;=\;
  \frac{\displaystyle\sum_{g \in \mathcal{G}_i} w_g\, N_g(h)}
       {\displaystyle\sum_{g \in \mathcal{G}_i^{\mathrm{sel}}} w_g\, D_g(h)},
  \qquad
  w_g = \frac{R_\mathrm{eff}(M_g)}{1 + z_g},
  \label{eq:lcat}
\end{equation}
% src: docs/derivations/G2c_gray_a9_a10_mapping.md §2, §4
where $M_g$ and $z_g$ are the catalogued massive-black-hole (MBH) host
mass and redshift of galaxy $g$, and the weight $w_g$ instantiates the
generic host weights $p(s\mid z)\,p(s\mid M)$ of \citet{Gray:2019ksv}
with the per-MBH EMRI rate $R_\mathrm{eff}(M)$ of \citet{Babak:2017tow}
and the detector-frame time dilation $1/(1+z)$.  The numerator sum
runs over the galaxies $\mathcal{G}_i$ inside the event's
three-dimensional localization region (galaxies outside it are
exponentially suppressed by equation~\eqref{eq:gwlike}); the choice of
the \emph{denominator} set $\mathcal{G}_i^{\mathrm{sel}}$ --- the same local
set, or the entire catalogue as in the published form of
\citet{Gray:2019ksv} --- is one of the two normalization choices at
the heart of this paper (Sections~\ref{sec:pitfall}
and~\ref{sec:estimators}).

The per-host numerator and selection factors are one-dimensional
redshift integrals against a host-redshift kernel $p_g(z)$,
\begin{align}
  N_g(h) &= \int \mathrm{d}z\;
  \mathcal{N}_3\!\big( (\phi_g, \theta_g, u(z,h));\; \hat\mu_i, \Sigma_i \big)\;
  p_g(z),
  \label{eq:ng}\\
  D_g(h) &= \int \mathrm{d}z\;
  p_\mathrm{det}\!\big(d_L(z,h), \Omega_g\big)\;
  p_g(z),
  \label{eq:dg}
\end{align}
% src: docs/derivations/G2c_gray_a9_a10_mapping.md §2 (N_g, D_g rows); derivations/dark_siren_likelihood.md Eq. (14.12)
evaluated over the redshift support of their integrands (the event's
$\pm 4\sigma_{d_L}$ distance image and the host's $\pm 4\sigma_z$
window, respectively), with $p_\mathrm{det}$ the detection probability
defined in Section~\ref{sec:framework:selection}.  The detection
probability appears in $D_g$ but never multiplies the measurement
Gaussian in $N_g$: selection effects are corrected in denominators
only \citep{Mandel:2018mve}.

The kernel $p_g(z)$ encodes what the catalogue's photometric redshift
$z_g \pm \sigma_{z,g}$ says about the host's true redshift.  We write
it in the general form
\begin{equation}
  p_g(z) \;=\;
  \frac{\mathcal{N}\big(z_g;\, z,\, \sigma_{z,g}^2\big)\;\Pi(z)}{Z_g},
  \qquad
  Z_g = \int \mathrm{d}z'\;
  \mathcal{N}\big(z_g;\, z',\, \sigma_{z,g}^2\big)\;\Pi(z'),
  \label{eq:kernel}
\end{equation}
% src: docs/derivations/G2b_host_z_volume_prior.md Eq. (1.4)
where $\mathcal{N}(z_g; z, \sigma_{z,g}^2)$ is the likelihood of the
catalogued value given a true redshift $z$, $\Pi(z)$ is a population
prior on host redshifts, and $Z_g$ normalizes each galaxy to unit
prior mass.  The bare choice $\Pi \equiv 1$ (kernel $=$ the catalogue
Gaussian itself) is the literal reading of the redshift distribution
$p(z_i)$ that \citet{Gray:2019ksv} treat as externally given; the
alternative, fixing $\Pi$ to the population prior of
equation~\eqref{eq:wpop} below, is the deconvolved kernel introduced
in Section~\ref{sec:estimators}.  For spectroscopic precision
($\sigma_z \to 0$) the kernel collapses to $\delta(z - z_g)$ and the
choice of $\Pi$ is irrelevant; at GLADE+ photometric widths it is
leading order (Section~\ref{sec:pitfall}).

An optional MBH-mass channel augments \eqref{eq:gwlike} with a fourth
coordinate, the redshifted-mass fraction $M_z/\hat M_z$ where
$M_z = M(1+z)$ and $M$ is the source-frame host MBH mass.  The
$M_z$ direction is marginalized analytically against a per-galaxy
host-mass kernel $p_g(M)$ using the standard conditional-Gaussian
decomposition \citep{Bishop2006},
and the corresponding selection factor integrates
$p_\mathrm{det}(d_L, M_z, \Omega_g)$ over the same kernels.  This
channel has no analogue in \citet{Gray:2019ksv}; its construction and
its rate-weighted mass prior are given in
Appendix~\ref{sec:app:eddm}.

\subsection{Completion term}
\label{sec:framework:completion}

Out-of-catalogue hosts are described by the population prior alone:
isotropic on the sky, $p(\Omega) = 1/4\pi$, and uniform in comoving
volume and source time, which in terms of the comoving volume element
per unit solid angle $\mathrm{d}V_\mathrm{c}/(\mathrm{d}z\,\mathrm{d}\Omega)$
reads
\begin{equation}
  w_\mathrm{pop}(z) \;\equiv\;
  \frac{1}{1+z}\,
  \frac{\mathrm{d}V_\mathrm{c}}{\mathrm{d}z\,\mathrm{d}\Omega}.
  \label{eq:wpop}
\end{equation}
% src: docs/derivations/G2b_host_z_volume_prior.md Eq. (1.3); docs/derivations/G2a_completion_sky_marginal_4pi.md Eq. (4)
Let $f(z, \Omega) \in [0,1]$ denote the (rate-weighted) completeness
of the catalogue, evaluated per HEALPix pixel, $f_k(z)$ for pixel $k$
\citep{Gray:2021sew}.
The numerator of the completion term is then the population-weighted,
incompleteness-weighted, sky-marginalized measurement likelihood
\begin{equation}
  B_\mathrm{num}(h) \;=\;
  \int_{z_-}^{z_+} \mathrm{d}z\;
  \big[1 - f_{k(\Omega_i)}(z)\big]\;
  \frac{\sin\hat\theta_i}{4\pi}\,
  \mathcal{N}\!\big(u(z,h);\, 1,\, \Sigma_{i,uu}\big)\;
  w_\mathrm{pop}(z),
  \label{eq:bnum}
\end{equation}
% src: docs/derivations/G2a_completion_sky_marginal_4pi.md §2 Eq. (5) and §3 Eq. (10)
where $[z_-, z_+]$ is the redshift image of the event's
$\pm 4\sigma_{d_L}$ distance support, $k(\Omega_i)$ is the pixel
containing the event's line of sight, and $\Sigma_{i,uu}$ is the
\emph{marginal} distance-fraction variance of $\Sigma_i$.  The factor
$(\sin\hat\theta_i / 4\pi)\,\mathcal{N}(u; 1, \Sigma_{i,uu})$ is the
exact narrow-beam marginal of the coordinate-density Gaussian
\eqref{eq:gwlike} over the isotropic sky prior --- not the Gaussian's
peak density; Section~\ref{sec:pitfall} shows that this distinction
alone changes the completion term by orders of magnitude, and
Appendix~\ref{sec:app:sky} derives the marginal and quantifies its
accuracy.  As in equation~\eqref{eq:ng}, no $p_\mathrm{det}$ factor
appears: the completion numerator is a likelihood numerator, and
selection lives in the denominators below.

\subsection{Selection normalization and the assembled likelihood}
\label{sec:framework:selection}

The detection probability $p_\mathrm{det}(d_L, \Omega)$ [and
$p_\mathrm{det}(d_L, M_z, \Omega)$ for the mass channel] is the
probability that an EMRI at that distance and sky position exceeds
$\rho_\mathrm{thr}$; it is calibrated from a dedicated injection
campaign run through the identical simulation pipeline, so that the
selection function of the inference matches that of the event
generator by construction.  Its integrals against the population
prior define the selection normalizations
\begin{align}
  D(h) &= \int_{z_\mathrm{min}}^{z_\mathrm{max}(h)} \mathrm{d}z\;
  \Big\langle p_\mathrm{det}\big(d_L(z,h), \Omega\big)
  \Big\rangle_{\Omega}\;
  w_\mathrm{pop}(z),
  \label{eq:Dh}\\
  \beta_{\bar G}(h) &= \int_{z_\mathrm{min}}^{z_\mathrm{max}(h)} \mathrm{d}z\;
  \Big\langle \big[1 - f(z,\Omega)\big]\,
  p_\mathrm{det}\big(d_L(z,h), \Omega\big) \Big\rangle_{\Omega}\;
  w_\mathrm{pop}(z),
  \label{eq:betabar}
\end{align}
% src: docs/derivations/G2c_gray_a9_a10_mapping.md §2 (D(h) and beta_Gbar rows)
where $\langle\cdot\rangle_\Omega$ is the isotropic sky average and
$z_\mathrm{max}(h)$ is the redshift horizon of the analysis
($p_\mathrm{det}$ vanishes beyond the injection horizon).  $D(h)$ is
the full-volume selection integral; $\beta_{\bar G}(h)$ its
out-of-catalogue restriction; and the \emph{catalogue selection
normalization}
\begin{equation}
  \beta_G(h) \;=\; D(h) - \beta_{\bar G}(h)
  \label{eq:betag}
\end{equation}
is the in-catalogue complement.  The hypothesis weights of
equation~\eqref{eq:mixture} are then
$p(G \mid D_i, h) = w_G(h) \equiv \beta_G(h)/D(h)$ and
$p(\bar G \mid D_i, h) = 1 - w_G(h) = \beta_{\bar G}(h)/D(h)$, and the
completion-term likelihood is
$p(x_i \mid \bar G, D_i, h) = B_\mathrm{num}(h)/\beta_{\bar G}(h)$.
Substituting into \eqref{eq:mixture}, the factors of
$\beta_{\bar G}$ cancel and the implemented per-event likelihood is
the single ratio
\begin{equation}
  p_i(h) \;=\;
  \frac{\beta_G(h)\, L_\mathrm{cat}(h) \;+\; B_\mathrm{num}(h)}{D(h)},
  \label{eq:assembled}
\end{equation}
% src: docs/derivations/G2c_gray_a9_a10_mapping.md §1
an exact algebraic rearrangement of \eqref{eq:mixture} that remains
finite in the complete-catalogue limit $f \to 1$
($\beta_{\bar G} \to 0$).  Both channels of the numerator carry the
same units ($\mathrm{Mpc^3\,sr^{-2}}$, from one population-volume
integral times one per-steradian GW density), so the two are
commensurable, and the residual constant left after dividing by
$D(h)$ is common to all $h$ and cancels when the posterior
\eqref{eq:posterior} is normalized.
% src: docs/derivations/G2a_completion_sky_marginal_4pi.md §6 (units of B_num and beta_G L_cat)

\subsection{Where prior consistency enters}
\label{sec:framework:consistency}

The population prior enters the likelihood in two guises.  It appears
\emph{explicitly}, as the measure $w_\mathrm{pop}(z)\,\mathrm{d}z$ of
every continuous integral --- the completion numerator
\eqref{eq:bnum} and the selection normalizations
\eqref{eq:Dh}--\eqref{eq:betabar}.  And it appears \emph{implicitly},
through whatever the in-catalogue kernel \eqref{eq:kernel} asserts
about where hosts live in redshift.  \emph{Prior consistency} is the
requirement that these two roles agree:
(i)~within each per-host ratio, the identical kernel $p_g(z)$ must
multiply the measurement factor in $N_g$
[equation~\eqref{eq:ng}] and the selection factor in $D_g$
[equation~\eqref{eq:dg}]; and
(ii)~across the catalogue and completion channels, every redshift
integral in the estimator must carry the population measure exactly
once --- never zero times, never twice.
% src: docs/derivations/G2b_host_z_volume_prior.md §3 ("dV_c counted once" rule)
Requirement (i) is satisfied by construction in
equation~\eqref{eq:lcat}.  Requirement (ii) constrains the pair
$\big(\Pi,\, \mathcal{G}_i^\mathrm{sel}\big)$: with the bare kernel
$\Pi \equiv 1$, the in-catalogue channel implicitly asserts a host
population that is uniform in $z$ while the completion and selection
integrals assert one uniform in comoving volume, and the mixture in
equation~\eqref{eq:assembled} weighs two mutually inconsistent priors
against each other through the $h$-dependent weight
$w_G(h) = \beta_G(h)/D(h)$ of equation~\eqref{eq:betag}.  The same requirement
applies verbatim to the mass channel: the host-mass kernel $p_g(M)$
must carry the rate prior $R_\mathrm{eff}(M)$ consistently between
its numerator and selection integrals (Appendix~\ref{sec:app:eddm}).
For narrow kernels the violation is second order, which is why it
escaped notice in spectroscopic-catalogue applications; at
photometric widths it is not.  Section~\ref{sec:pitfall} demonstrates
that this seemingly formal bookkeeping choice is the difference
between a posterior railed at the prior boundary and a calibrated
one.
